\lyseis{4}{Ευκλείδεια διαίρεση: το υπόλοιπο που μένει}
\sol{1} α) $4.002:69=58$ (δοκιμή: $69\cdot 58=4.002$) \quad β) $1.445:17=85$ \quad γ) $925:37=25$ \quad δ) $3.621:213=17$ \quad ε) $35.280:2.940=12$ \quad στ) $5.082:77=66$. Όλες είναι τέλειες διαιρέσεις, με υπόλοιπο 0.
\sol{2} α) $43=7\cdot 6+1$ \quad β) $100=9\cdot 11+1$ \quad γ) $245=12\cdot 20+5$ \quad δ) $1.000=33\cdot 30+10$ \quad ε) $59=60\cdot 0+59$ (το 60 δεν χωράει στο 59, πηλίκο 0) \quad στ) $0=5\cdot 0+0$.
\sol{3} α) Ναι, με διαιρέτη το 35 ($\upsilon=20<35$, πηλίκο 3). Όχι με διαιρέτη το 3, γιατί $20>3$. β) Όχι: το υπόλοιπο 40 είναι μεγαλύτερο και από το 38 και από το 19. γ) Ναι, με διαιρέτη το 42 (πηλίκο 35) και με διαιρέτη το 35 (πηλίκο 42), γιατί $30<35$. δ) Ναι, με διαιρέτη το 18 (πηλίκο 16) και με διαιρέτη το 16 (πηλίκο 18), γιατί $12<16$.
\sol{4} Τα υπόλοιπα είναι μικρότερα από το 8: 0, 1, 2, 3, 4, 5, 6 ή 7.
\sol{5} $\Delta=9\cdot 73+4=657+4=661$. Δοκιμή: $661:9$ δίνει πηλίκο 73 και υπόλοιπο 4.
\sol{6} α) $\Delta=12\cdot 7+5=89$ \quad β) $100=8\cdot 12+4$, άρα $\pi=12$, $\upsilon=4$ \quad γ) $59=\delta\cdot 4+3$, άρα $4\delta=56$, $\delta=14$. Έλεγχος: $3<14$.
\sol{7} α) Λάθος: το υπόλοιπο είναι μικρότερο από τον διαιρέτη, άρα το πολύ 5. β) Σωστό. γ) Λάθος: δεν διαιρούμε με το 0. δ) Σωστό. ε) Λάθος: $17:1=17$. στ) Σωστό.
\sol{8} Η ισότητα ισχύει: $12\cdot 6+13=72+13=85$. Όμως το υπόλοιπο 13 δεν είναι μικρότερο από τον διαιρέτη 12: ο 12 χωράει άλλη μία φορά στο 13. Σωστό: $85=12\cdot 7+1$.
\sol{9} $387=49\cdot 7+44$: χρειάζονται 8 λεωφορεία. Στο όγδοο υπάρχουν 44 παιδιά και μένουν $49-44=5$ κενές θέσεις. Με 7 λεωφορεία μένουν στο σχολείο 44 παιδιά.
\sol{10} $150=7\cdot 21+3$: Τρίτη και 3 μέρες, δηλαδή Παρασκευή. Πριν από 8 ημέρες: $8=7+1$, άρα μία μέρα πριν από την Τρίτη, δηλαδή Δευτέρα.
\sol{11} $1.000=33\cdot 30+10$: γεμίζει 30 γλάστρες και περισσεύουν 10 σπόροι. Όχι, δεν φτάνουν για άλλη γλάστρα, γιατί $10<33$.
\sol{12} $87=13\cdot 6+9$: αγοράζει 6 μπάλες (78 €) και του μένουν 9 €.
\sol{13} $168:3=56$, $168:4=42$, $168:6=28$, $168:7=24$: τέλειες διαιρέσεις, μπορούν να σταθούν σε ίσες ομάδες. Το $168:5$ δίνει πηλίκο 33 και υπόλοιπο 3 ($168=5\cdot 33+3$): περισσεύουν 3 παιδιά.
\sol{14} $50=24\cdot 2+2$: μετά από 2 ολόκληρα εικοσιτετράωρα το ρολόι δείχνει πάλι 10:00 και μένουν 2 ώρες. Θα είναι 12:00, δύο μέρες αργότερα.
\sol{15} Οι αριθμοί είναι 24, 31, 38 ($24=7\cdot 3+3$, $31=7\cdot 4+3$, $38=7\cdot 5+3$). Η διαφορά ανάμεσα σε δύο διαδοχικούς είναι πάντα 7, δηλαδή ο διαιρέτης.
\sol{16} Για $\nu=10$: $10:3$ έχει υπόλοιπο 1. Για 16: $16=3\cdot 5+1$, υπόλοιπο 1. Για 22: $22=3\cdot 7+1$, υπόλοιπο 1. Πάντα 1. Ο λόγος: $\nu=6\cdot\kappa+4$, το $6\kappa$ διαιρείται με το 3, και το 4 αφήνει υπόλοιπο 1.
\sol{17} Ο διαιρέτης πρέπει να είναι μεγαλύτερος από το υπόλοιπο 7, άρα ο μικρότερος είναι ο 8. Τότε $\Delta=8\cdot 12+7=96+7=103$.
\sol{18} Αν $6:0=x$, τότε $6=0\cdot x$, που είναι αδύνατο, γιατί το $0\cdot x$ είναι πάντα 0. Δεν υπάρχει $x$. Στο $0:0=x$ θα έπρεπε $0=0\cdot x$, που ισχύει για κάθε $x$: δεν υπάρχει ένας μοναδικός αριθμός. Και στις δύο περιπτώσεις η πράξη δεν ορίζεται.
\sol{19} $1.000.000=7\cdot 142.857+1$. Το υπόλοιπο είναι 1, άρα μία μέρα μετά την Τρίτη: Τετάρτη.
